\documentclass[11pt]{article}
\usepackage[a4paper,margin=1.1in]{geometry}
\usepackage{amsmath,amssymb,amsthm,mathtools}
\usepackage{enumitem}
\usepackage{hyperref}
\usepackage{xcolor}
\usepackage{microtype}
\usepackage{booktabs}

\hypersetup{colorlinks=true,linkcolor=blue!60!black,citecolor=blue!60!black,urlcolor=blue!60!black}

\newtheorem{theorem}{Theorem}[section]
\newtheorem{lemma}[theorem]{Lemma}
\newtheorem{proposition}[theorem]{Proposition}
\newtheorem{corollary}[theorem]{Corollary}
\newtheorem{conjecture}[theorem]{Conjecture}
\newtheorem{question}[theorem]{Question}
\theoremstyle{definition}
\newtheorem{definition}[theorem]{Definition}
\newtheorem{remark}[theorem]{Remark}
\newtheorem{example}[theorem]{Example}

\newcommand{\Traceace}{\operatorname{Tr}}
\newcommand{\one}{\mathbf 1}
\newcommand{\E}{\mathbb E}
\newcommand{\R}{\mathbb R}
\newcommand{\eps}{\varepsilon}
\newcommand{\same}{\mathrm{same}}
\newcommand{\cross}{\mathrm{cross}}
\newcommand{\dd}{\,d}
\newcommand{\ip}[2]{\langle #1,#2\rangle}
\newcommand{\norm}[1]{\lVert #1\rVert}

\title{A Flag-Algebra and Stability Framework for the Odd-Cycle Graphon Inequality}
\author{Working note generated for the odd-cycle project}
\date{\today}

\begin{document}
\maketitle

\begin{abstract}
We record a self-contained formulation of the flag-algebra/stability attack on the conjectural inequality
\[
        t(C_{2k+1},W)\ge p^{2k+1}-p(1-p)^{2k},
        \qquad p=t(K_2,W),
\]
for graphons.  The note does not claim a proof of the general conjecture.  Instead it explains the exact quantum-graph reformulation, the first- and second-variation calculations at all balanced Turan graphons, the type-2 rooted path-kernel square that captures the Turan KKT equations, and the resulting stability and certificate-search targets.  The goal is to give collaborators a precise and reusable starting point for a flag-algebra/stability proof.
\end{abstract}

\tableofcontents

\section{The conjecture and notation}

A \emph{graphon} is a symmetric measurable function
\[
        W:[0,1]^2\to[0,1].
\]
For a finite graph $F$, its homomorphism density in $W$ is
\[
        t(F,W)=\int_{[0,1]^{V(F)}}\prod_{ij\in E(F)} W(x_i,x_j)\prod_{i\in V(F)}dx_i.
\]
Let
\[
        p=t(K_2,W)=\int_{[0,1]^2}W(x,y)\dd x\dd y,
        \qquad q=1-p.
\]
Our main conjecture is the following.

\begin{conjecture}[Odd-cycle lower bound]
For every graphon $W$ and every $k\ge1$,
\[
        t(C_{2k+1},W)\ge p^{2k+1}-p(1-p)^{2k}.
\]
Equivalently, writing $m=2k+1$,
\[
        t(C_m,W)\ge p^m-pq^{m-1}.
\]
\end{conjecture}

The right-hand side is nonpositive when $p\le1/2$, so the difficult regime is
\[
        p>1/2.
\]
The conjecture is sharp at balanced complete multipartite graphons.  If $T_r$ is the balanced complete $r$-partite graphon, then
\[
        p=t(K_2,T_r)=1-\frac1r,
\]
and
\[
        t(C_m,T_r)=p^m-p(1-p)^{m-1}
\]
for every odd $m$.  Indeed, $t(C_m,T_r)$ is the weighted probability that a random map from $C_m$ to $[r]$ is a proper coloring, and the chromatic polynomial of an odd cycle is
\[
        \chi(C_m,r)=(r-1)^m-(r-1).
\]

\begin{remark}[What this note is and is not]
Earlier parts of the project established the conjecture for several special regimes, including all regular graphons, $C_3,C_5,C_7,C_9,C_{11}$, and several all-$k$ structural families.  This note focuses on the subsequent flag-algebra/stability direction.  It does not contain a global proof.  Its purpose is to state the exact objects that a flag-algebra certificate or stability proof should use.
\end{remark}

\section{Quantum-graph reformulation}

For fixed odd $m=2k+1$, define the defect functional
\[
        \Phi_m(W)=t(C_m,W)-p^m+p(1-p)^{m-1}.
\]
The conjecture is
\[
        \Phi_m(W)\ge0.
\]

Because $m-1$ is even,
\[
        p(1-p)^{m-1}=p\sum_{j=0}^{m-1}(-1)^j\binom{m-1}{j}p^j
        =\sum_{s=1}^{m}(-1)^{s-1}\binom{m-1}{s-1}p^s,
\]
and the coefficient of $p^m$ is $+1$.  Therefore the $-p^m$ term cancels, leaving
\[
        \Phi_m(W)
        =t(C_m,W)+\sum_{s=1}^{m-1}(-1)^{s-1}\binom{m-1}{s-1}p^s.
\]
Since $p^s=t(sK_2,W)$, where $sK_2$ denotes the disjoint union of $s$ edges, we obtain the linear quantum-graph form.

\begin{proposition}[Quantum-graph form]
Let
\[
        Q_m=C_m+
        \sum_{s=1}^{m-1}(-1)^{s-1}\binom{m-1}{s-1}sK_2.
\]
Then
\[
        \Phi_m(W)=t(Q_m,W).
\]
Consequently, the conjecture for a fixed odd $m$ is exactly the quantum-graph inequality
\[
        t(Q_m,W)\ge0
\]
for all graphons $W$.
\end{proposition}

\begin{proof}
The computation above gives the identity, using multiplicativity of homomorphism densities over disjoint unions:
\[
        t(sK_2,W)=t(K_2,W)^s=p^s.
\]
\end{proof}

This is an important point.  For each fixed $m$, the conjecture is not a nonlinear graphon inequality after introducing disjoint-edge quantum graphs.  It is a linear inequality in the algebra of quantum graphs, exactly the kind of object accessible to flag algebras.

\section{The rooted path kernel}

Let $T_W$ be the self-adjoint Hilbert--Schmidt operator
\[
        T_W f(x)=\int_0^1 W(x,y)f(y)\dd y.
\]
For $n\ge1$, write $K_n^W$ for the kernel of $T_W^n$.  Thus
\[
        K_n^W(x,y)=\int_{[0,1]^{n-1}} W(x,z_1)W(z_1,z_2)\cdots W(z_{n-1},y)\dd z_1\cdots \dd z_{n-1}.
\]
This is the density of an $n$-edge path rooted at its endpoints $x,y$.

For $m=2k+1$, set
\[
        n=m-1=2k.
\]
Then
\[
        t(C_m,W)=\int_{[0,1]^2} W(x,y)K_n^W(x,y)\dd x\dd y.
\]
Define
\[
        A_m(p)=p^n-q^n,
        \qquad
        B_m(p)=p^n+p q^{n-1}.
\]
These constants are chosen because at Turan graphons the rooted path kernel takes precisely these two values.

\begin{lemma}[Turan path-kernel values]
Let $T_r$ be the balanced complete $r$-partite graphon, and put
\[
        p_r=1-\frac1r,
        \qquad q_r=\frac1r.
\]
Then
\[
        K_n^{T_r}(x,y)=A_m(p_r)
        \quad\text{for cross-part pairs,}
\]
and
\[
        K_n^{T_r}(x,y)=B_m(p_r)
        \quad\text{for same-part pairs.}
\]
\end{lemma}

\begin{proof}
The nonzero spectrum of $T_r$ consists of $p_r$ on the constants and $-q_r$ on the $(r-1)$-dimensional mean-zero space of step functions that are constant on the parts.  Since $n$ is even,
\[
        T_r^n=p_r^n J+q_r^n P_0,
\]
where $J$ is the projection onto constants and $P_0$ is the projection onto that mean-zero step space.  The kernel of $P_0$ is
\[
        r\mathbf 1_{\same}(x,y)-1.
\]
Thus
\[
        K_n^{T_r}(x,y)=p_r^n+q_r^n(r\mathbf 1_{\same}(x,y)-1).
\]
On cross-pairs this is $p_r^n-q_r^n=A_m(p_r)$.  On same-part pairs this is
\[
        p_r^n+(r-1)q_r^n=p_r^n+p_rq_r^{n-1}=B_m(p_r).
\]
\end{proof}

The two constants satisfy the useful identity
\[
        pA_m(p)+qB_m(p)=p^n.
\]
Indeed,
\[
        p(p^n-q^n)+q(p^n+pq^{n-1})=p^n.
\]

The defect has the rooted-kernel expression
\[
        \Phi_m(W)=\int W(x,y)(K_n^W(x,y)-A_m(p))\dd x\dd y.
\]
Indeed,
\[
        \int W(K_n-A_m)=t(C_m,W)-pA_m(p)=t(C_m,W)-p^{m}+pq^{m-1}.
\]

\section{First variation and the Turan KKT multipliers}

Let
\[
        \varphi_m(p)=p^m-pq^{m-1},
        \qquad q=1-p.
\]
Then
\[
        \Phi_m(W)=t(C_m,W)-\varphi_m(p).
\]
We compute the first variation at the balanced Turan graphons.

\begin{lemma}[First variation of cycle density]
For any bounded symmetric perturbation $H$,
\[
        D t(C_m,W)[H]
        =m\int H(x,y)K_{m-1}^W(x,y)\dd x\dd y.
\]
Also,
\[
        Dp[H]=\int H(x,y)\dd x\dd y.
\]
\end{lemma}

\begin{proof}
Differentiate the product defining $t(C_m,W)$ edge-by-edge.  Cyclic symmetry gives $m$ identical terms, each equal to the displayed integral.  The formula for $Dp$ is immediate.
\end{proof}

A direct derivative gives
\[
        \varphi_m'(p)=mp^{m-1}-q^{m-1}+(m-1)pq^{m-2}.
\]
Since $n=m-1$, it is convenient to define
\[
        \Theta_m(p)=\frac1m\varphi_m'(p).
\]
Then
\[
        \Theta_m(p)=A_m(p)+\frac{m-1}{m}q^{m-2}
\]
and
\[
        \Theta_m(p)=B_m(p)-\frac1m q^{m-2}.
\]
Thus $\Theta_m(p)$ lies strictly between the Turan edge and nonedge rooted path-kernel values whenever $0<q<1$.

\begin{theorem}[Turan KKT first-variation formula]
Let $T_r$ be the balanced complete $r$-partite graphon.  Let $H$ be a bounded symmetric perturbation such that $T_r+\eps H$ is graphon-valued for all sufficiently small $\eps>0$.  Equivalently,
\[
        H\ge0\quad\text{inside the diagonal zero-blocks,}
        \qquad
        H\le0\quad\text{on the cross one-blocks.}
\]
Then
\[
        \left.\frac{d}{d\eps}\right|_{\eps=0+}\Phi_m(T_r+\eps H)
        =q_r^{m-2}\left(\int_{\same}H-(m-1)\int_{\cross}H\right),
\]
where $q_r=1/r$.
Consequently the one-sided derivative is nonnegative, and it is positive unless $H=0$ almost everywhere on the support of the perturbation.
\end{theorem}

\begin{proof}
Using the previous lemma,
\[
        D\Phi_m(T_r)[H]
        =m\int H K_{m-1}^{T_r}-\varphi_m'(p_r)\int H.
\]
By the Turan path-kernel calculation,
\[
        K_{m-1}^{T_r}=p_r^{m-1}+q_r^{m-1}(r\mathbf 1_{\same}-1).
\]
Thus on same-part pairs,
\[
        K_{\same}=p_r^{m-1}+p_rq_r^{m-2},
\]
while on cross-pairs,
\[
        K_{\cross}=p_r^{m-1}-q_r^{m-1}.
\]
The coefficient multiplying $H$ on same-part pairs is
\[
\begin{aligned}
        mK_{\same}-\varphi_m'(p_r)
        &=m(p_r^{m-1}+p_rq_r^{m-2})
          -\bigl(mp_r^{m-1}-q_r^{m-1}+(m-1)p_rq_r^{m-2}\bigr)\\
        &=p_rq_r^{m-2}+q_r^{m-1}=q_r^{m-2}.
\end{aligned}
\]
The coefficient on cross-pairs is
\[
\begin{aligned}
        mK_{\cross}-\varphi_m'(p_r)
        &=m(p_r^{m-1}-q_r^{m-1})
          -\bigl(mp_r^{m-1}-q_r^{m-1}+(m-1)p_rq_r^{m-2}\bigr)\\
        &=-(m-1)q_r^{m-1}-(m-1)p_rq_r^{m-2}\\
        &=-(m-1)q_r^{m-2}.
\end{aligned}
\]
The displayed formula follows.  The sign assertion follows from the feasibility conditions on $H$.
\end{proof}

\begin{remark}
This calculation is the KKT picture behind the conjecture.  On zero-blocks, increasing $W$ has a positive first-order cost.  On one-blocks, decreasing $W$ has a positive first-order cost because $H\le0$ and the multiplier is negative.  Hence balanced Turan graphons are one-sided local minimizers in all edge-value directions on the fixed partition.
\end{remark}

\section{Second variation in part-size directions}

The previous section fixes the balanced partition and perturbs edge values.  We also need to understand perturbations of the part sizes.

Let $a_1,\ldots,a_r$ be positive numbers summing to $1$, and let $T(a_1,\ldots,a_r)$ be the complete $r$-partite graphon with part sizes $a_i$.  Put
\[
        a_i=\frac1r+\eps_i,
        \qquad
        \sum_i\eps_i=0.
\]
Then
\[
        p(a)=t(K_2,T(a))=1-\sum_i a_i^2=p_r-\sum_i\eps_i^2.
\]
The following calculation records the second-order stability of the balanced part sizes.

\begin{theorem}[Second-order part-size stability]
For fixed odd $m\ge3$ and fixed $r\ge2$,
\[
        \Phi_m(T(a_1,\ldots,a_r))
        =\Gamma_{m,r}\sum_{i=1}^r\eps_i^2
        +O_{m,r}\left(\sum_i |\eps_i|^3\right),
\]
where
\[
        \Gamma_{m,r}=\frac{P_m(r-1)}{r^{m-2}}>0,
\]
and, writing $R=r-1$,
\[
        P_m(R)=
        m\frac{R^{m-1}-(m-2)R^2-(m-3)R}{(R+1)^2}
        +\frac{m(m-1)}2-1.
\]
\end{theorem}

\begin{proof}
The homomorphism density $t(C_m,T(a))$ is the weighted partition function of proper colorings of $C_m$ with colors $1,\ldots,r$ and color weights $a_i$:
\[
        t(C_m,T(a))=\sum_{i_1,\ldots,i_m\,:\, i_j\ne i_{j+1}} a_{i_1}\cdots a_{i_m},
\]
where $i_{m+1}=i_1$.  Equivalently it is the trace of the $m$-th power of the $r\times r$ matrix
\[
        M(a)_{ij}=\begin{cases}
        \sqrt{a_i a_j},& i\ne j,\\
        0,& i=j.
        \end{cases}
\]
At $a_i=1/r$, the matrix has eigenvalues $p_r=R/(R+1)$ and $-1/r=-1/(R+1)$ with multiplicity $r-1$.

Expanding the weighted coloring partition function to second order in $\eps_i$ yields
\[
        t(C_m,T(a))=t(C_m,T_r)+C_{m,r}\sum_i\eps_i^2+O_{m,r}\left(\sum_i|\eps_i|^3\right),
\]
where
\[
        C_{m,r}=\frac{m}{r^{m-2}}
        \frac{R^{m-1}-(m-2)R^2-(m-3)R}{(R+1)^2}.
\]
On the other hand,
\[
        p(a)=p_r-\sum_i\eps_i^2,
\]
so
\[
        \varphi_m(p(a))=\varphi_m(p_r)-\varphi_m'(p_r)\sum_i\eps_i^2+O\left((\sum_i\eps_i^2)^2\right).
\]
A direct substitution gives
\[
        \varphi_m'(p_r)=\frac{1}{r^{m-2}}\left(\frac{m(m-1)}2-1\right).
\]
Therefore the coefficient of $\sum_i\eps_i^2$ in
\[
        \Phi_m=t(C_m,T(a))-\varphi_m(p(a))
\]
is exactly $\Gamma_{m,r}=P_m(R)/r^{m-2}$.

It remains to prove positivity.  We have
\[
        P_m(1)=m-1>0,
\]
and differentiating gives
\[
        P_m'(R)=
        \frac{m\bigl((m-3)(R^{m-1}-1)+(m-1)(R^{m-2}-R)\bigr)}{(R+1)^3}.
\]
For odd $m\ge3$ and $R\ge1$, both $R^{m-1}-1$ and $R^{m-2}-R$ are nonnegative.  Thus $P_m'(R)\ge0$ for $R\ge1$, and so $P_m(R)\ge P_m(1)>0$.
\end{proof}

\begin{remark}
The explicit second-order coefficient is a diagnostic for stability.  It says that balanced Turan graphons are strict local minimizers not only with respect to edge-value perturbations on a fixed partition, but also with respect to nontrivial part-size perturbations within the complete $r$-partite family.
\end{remark}

\section{The type-2 KKT square residual}

The first variation suggests a type-2 flag square.  Recall
\[
        A_m(p)=p^n-q^n,
        \qquad
        B_m(p)=p^n+pq^{n-1},
        \qquad n=m-1.
\]
Define
\[
\boxed{
        \mathcal S_m(W)=
        \int W(x,y)(K_n^W(x,y)-A_m(p))^2\dd x\dd y
        +\int (1-W(x,y))(K_n^W(x,y)-B_m(p))^2\dd x\dd y.
        }
\]
This is a nonnegative quantity.  It vanishes at every balanced Turan graphon $T_r$, because $K_n^{T_r}=A_m(p_r)$ on edges and $K_n^{T_r}=B_m(p_r)$ on nonedges.

Let
\[
        R_m^W(x,y)=mK_n^W(x,y)-\varphi_m'(p).
\]
Then, since $\Theta_m(p)=\varphi_m'(p)/m$,
\[
        R_m^W+(m-1)q^{n-1}=m(K_n^W-A_m(p)),
\]
and
\[
        R_m^W-q^{n-1}=m(K_n^W-B_m(p)).
\]
Therefore
\[
\boxed{
        m^2\mathcal S_m(W)=
        \int W\bigl(R_m^W+(m-1)q^{n-1}\bigr)^2
        +\int(1-W)\bigl(R_m^W-q^{n-1}\bigr)^2.
        }
\]
This is the precise square of the deviation from the Turan KKT multipliers.

\subsection{Expansion as ordinary graph densities}

The square residual is a flag object, but it can also be written as an ordinary quantum graph.  Let $\Theta_{n,n,1}$ denote the theta graph consisting of two internally disjoint $n$-edge paths between two common endpoints, plus the edge between the endpoints.  Then
\[
        \int W(K_n^W)^2=t(\Theta_{n,n,1},W).
\]
Also,
\[
        \int (K_n^W)^2=t(C_{2n},W),
\]
and
\[
        \int K_n^W=t(P_n,W).
\]
Expanding the square gives
\[
\begin{aligned}
        \mathcal S_m(W)
        ={}&t(C_{2n},W)
        +2(B_m(p)-A_m(p))t(C_m,W)\\
        &-2B_m(p)t(P_n,W)
        +A_m(p)^2p+B_m(p)^2q.
\end{aligned}
\]
Since $A_m(p)$ and $B_m(p)$ are polynomials in $p$, this can be expanded into a linear combination of ordinary graph densities involving $C_{2n}$, $C_m$, $P_n$, and disjoint unions with edges.  In flag-algebra language, it is a type-2 sum of squares.

\section{A forcing lemma for counterexamples}

Let
\[
        \Delta_m(W)=\Phi_m(W)=t(C_m,W)-\bigl(p^m-pq^{m-1}\bigr).
\]
Then
\[
        \Delta_m(W)=\int W(K_n^W-A_m(p)).
\]
Define the nonedge companion
\[
        \Xi_m(W)=\int (1-W)(K_n^W-B_m(p)).
\]
Using $pA_m+qB_m=p^n$, we obtain
\[
\begin{aligned}
        \Delta_m(W)+\Xi_m(W)
        &=\int K_n^W-p^n\\
        &=t(P_n,W)-p^n.
\end{aligned}
\]
By the Blakley--Roy inequality for paths, or equivalently Sidorenko's theorem for paths,
\[
        t(P_n,W)\ge p^n.
\]
Hence
\[
        \Delta_m(W)+\Xi_m(W)\ge0.
\]

\begin{proposition}[Counterexamples force square mass]
If $W$ is a counterexample with
\[
        \Delta_m(W)=-\delta<0,
\]
then
\[
        \mathcal S_m(W)\ge \delta^2\left(\frac1p+\frac1q\right)=\frac{\delta^2}{pq}.
\]
Equivalently,
\[
        \Delta_m(W)\ge -\sqrt{pq\,\mathcal S_m(W)}.
\]
\end{proposition}

\begin{proof}
If $\Delta_m(W)=-\delta$, then the identity above gives $\Xi_m(W)\ge\delta$.  By Cauchy--Schwarz,
\[
        \int W(K_n-A_m)^2\ge \frac{\delta^2}{\int W}=\frac{\delta^2}{p},
\]
and
\[
        \int (1-W)(K_n-B_m)^2\ge \frac{\delta^2}{\int(1-W)}=\frac{\delta^2}{q}.
\]
Adding proves the claim.
\end{proof}

\begin{remark}
This is not a proof of the conjecture.  It says that any negative counterexample must be quantitatively far from the Turan KKT square-zero locus.  Thus the type-2 square residual is not cosmetic; it is forced by any attempted counterexample.
\end{remark}

\section{KKT conditions for a hypothetical minimizer}

Suppose the conjecture is false for a fixed odd $m$.  By compactness of graphon space modulo weak isomorphism and continuity of homomorphism densities, there is a minimizer $W_*$ of $\Phi_m$ with
\[
        \Phi_m(W_*)<0.
\]
Let
\[
        p_*=t(K_2,W_*),
        \qquad
        K_*=K_n^{W_*},
\]
and
\[
        R_*(x,y)=mK_*(x,y)-\varphi_m'(p_*).
\]
The box constraints $0\le W\le1$ imply the following KKT conditions.

\begin{proposition}[Threshold condition for an interior minimizer]
At a minimizing counterexample $W_*$,
\[
        R_*(x,y)\ge0\quad\text{where } W_*(x,y)=0,
\]
\[
        R_*(x,y)\le0\quad\text{where } W_*(x,y)=1,
\]
and
\[
        R_*(x,y)=0\quad\text{where }0<W_*(x,y)<1.
\]
Equivalently, writing $\Theta_m(p_*)=\varphi_m'(p_*)/m$,
\[
        K_*(x,y)\ge\Theta_m(p_*)\quad\text{on zero-pairs,}
\]
\[
        K_*(x,y)\le\Theta_m(p_*)\quad\text{on one-pairs,}
\]
and equality holds on fractional pairs.
\end{proposition}

\begin{proof}
Consider a feasible perturbation that increases $W$ on a set where $W_*=0$.  Minimality forces the one-sided derivative to be nonnegative there, giving $R_*\ge0$.  Similarly, decreasing $W$ on a set where $W_*=1$ must not decrease the functional, giving $R_*\le0$.  On a fractional region both signs of perturbation are feasible, hence $R_*=0$.
\end{proof}

Thus any genuine counterexample must be a self-consistent threshold graphon of its own even-path kernel.  Balanced Turan graphons satisfy the same threshold relation, with
\[
        A_m(p)<\Theta_m(p)<B_m(p),
\]
and the gaps
\[
        \Theta_m(p)-A_m(p)=\frac{m-1}{m}q^{m-2},
        \qquad
        B_m(p)-\Theta_m(p)=\frac1m q^{m-2}.
\]

\section{Local stability around balanced Turan graphons}

The variation formulae yield a local stability result in a standard form.

\begin{proposition}[Fixed-partition linear stability]
Fix $m$ and $r$.  There exists $\eps_0=\eps_0(m,r)>0$ such that the following holds.  Let $W=T_r+H$ on the balanced $r$-partition, with
\[
        \norm{H}_\infty\le\eps_0,
\]
and with the graphon constraints forcing
\[
        H\ge0\text{ on same-part blocks},
        \qquad
        H\le0\text{ on cross blocks}.
\]
Then
\[
        \Phi_m(W)
        \ge\frac12q_r^{m-2}\left(
        \int_{\same}H+\int_{\cross}(-H)
        \right).
\]
\end{proposition}

\begin{proof}
By Taylor expansion,
\[
        \Phi_m(T_r+H)=\Phi_m(T_r)+D\Phi_m(T_r)[H]+O_{m,r}(\norm{H}_\infty \norm{H}_1),
\]
where the remainder follows because $\Phi_m$ is a polynomial functional of degree $m$ in $W$ and degree $m$ in $p$.  Since $\Phi_m(T_r)=0$ and
\[
        D\Phi_m(T_r)[H]
        \ge q_r^{m-2}\left(\int_{\same}H+\int_{\cross}(-H)\right),
\]
choosing $\eps_0$ small enough gives the result.
\end{proof}

Together with second-order part-size stability, this gives a robust local picture: balanced Turan graphons are strict local minimizers in the natural graphon directions that preserve proximity to the Turan template.

\section{How this suggests a flag-algebra certificate}

The calculations above suggest that a fixed-$m$ flag-algebra certificate should not be sought as an arbitrary SOS identity.  It should be organized around the KKT square $\mathcal S_m$.

A plausible certificate shape is
\[
        \Phi_m(W)+\lambda_m(p)\mathcal S_m(W)+
        \sum_d \mu_{m,d}(p)\mathcal K_{m,d}(W)\ge0,
\]
where:
\begin{itemize}[leftmargin=2em]
        \item $\mathcal S_m$ is the type-2 KKT square residual;
        \item the $\mathcal K_{m,d}$ are nonnegative path-layer or Schur-convexity surplus terms, for example the even-layer terms appearing in the Kim--Lee extended commonality theorem;
        \item $\lambda_m$ and $\mu_{m,d}$ are nonnegative functions of $p,q$;
        \item the remaining expression should vanish at every balanced Turan graphon and factor through Turan-stability terms.
\end{itemize}

The reason $\mathcal S_m$ is essential is the forcing lemma: a counterexample must have square mass at least $\delta^2/(pq)$.  Thus any proof that rules out counterexamples should naturally use this mass.

\subsection{Type-2 flags}

The natural type has two labelled vertices.  The two basic type-2 flags are:
\begin{itemize}[leftmargin=2em]
        \item the labelled pair joined by an edge;
        \item the labelled pair not joined by an edge.
\end{itemize}
The rooted path kernel $K_n^W(x,y)$ is the density of an $n$-edge path between the labelled vertices.  Therefore
\[
        W(x,y)(K_n^W-A_m(p))^2
\]
and
\[
        (1-W(x,y))(K_n^W-B_m(p))^2
\]
are exactly squares of type-2 flag expressions, multiplied by the edge or nonedge flag.

A computational flag-algebra search should include these squares explicitly as distinguished blocks rather than hoping that a generic SDP discovers them.

\section{Relation to the complement path-cycle identity}

Previous work in the project used the complement $U=1-W$ and exact inclusion--exclusion.  For $m=2k+1$ and $n=m-1$, let
\[
        x_n=t(P_n,U),
        \qquad
        c_m=t(C_m,U).
\]
There is an exact identity of the form
\[
        \Delta_m(W)=\Pi_m(U)+(x_n-c_m),
\]
where $\Pi_m(U)$ is the complement path-certificate expression obtained after expanding $t(C_m,1-U)$ and comparing with $(1-q)^m-(1-q)q^{m-1}$.

The crude path-certificate method tries to prove $\Pi_m(U)\ge0$, using only $x_n-c_m\ge0$.  This works for small cycles and in partial ranges, but it fails from larger $m$ onward.  The flag-square method explains why: the discarded correction $x_n-c_m$ contains rooted two-endpoint information.  In fact,
\[
        x_{2a}-c_{2a+1}=\int (1-U(x,y))\bigl(K_a^U(x,y)\bigr)^2\dd x\dd y.
\]
For $C_{13}$ this gives
\[
        x_{12}-c_{13}=\int(1-U)(K_6^U)^2.
\]
Thus the missing correction is itself a type-2 square.  This is consistent with the general $\mathcal S_m$ philosophy.

\section{Concrete SDP/search target}

For a fixed odd $m$, the next computational target is:
\begin{enumerate}[leftmargin=2em]
        \item Construct the quantum graph
        \[
                Q_m=C_m+
                \sum_{s=1}^{m-1}(-1)^{s-1}\binom{m-1}{s-1}sK_2.
        \]
        \item Work in the flag algebra of type $2$ with flags up to a size large enough to contain the rooted $n$-path terms, where $n=m-1$.
        \item Include the KKT square residual $\mathcal S_m$ as an explicit SOS block.
        \item Add known nonnegative even path-layer surplus terms, especially those coming from the Schur-convexity mechanism behind extended commonality for paths and cycles.
        \item Impose vanishing at all balanced Turan graphons.  Equivalently, enforce the relations
        \[
                K_n=A_m(p)\quad\text{on edges},
                \qquad
                K_n=B_m(p)\quad\text{on nonedges}
        \]
        at the Turan templates.
\end{enumerate}

A successful certificate should reflect the KKT multipliers from the first variation:
\[
        q_r^{m-2}\quad\text{on zero-blocks},
        \qquad
        -(m-1)q_r^{m-2}\quad\text{on one-blocks}.
\]
This is a strong diagnostic constraint on any SDP output.

\section{What is proved, and what remains}

The results in this note establish:
\begin{itemize}[leftmargin=2em]
        \item the exact quantum-graph form of the conjecture;
        \item equality at all balanced Turan graphons;
        \item the exact first-variation/KKT multipliers at balanced Turan graphons;
        \item strict second-order stability in part-size directions;
        \item the type-2 rooted path-kernel square residual $\mathcal S_m$;
        \item a forcing lemma showing that every counterexample must have nontrivial $\mathcal S_m$ mass;
        \item KKT threshold conditions for any hypothetical global counterexample;
        \item a precise flag-algebra SDP ansatz.
\end{itemize}

What remains is the hard global step:
\[
        \text{prove }t(Q_m,W)\ge0\text{ globally for every graphon }W.
\]
The present framework suggests that this should be approached by a flag-algebra certificate plus stability, not by discarding the path-cycle correction or by one-variable spectral estimates alone.

\section{Suggested tasks for collaborators}

\begin{enumerate}[leftmargin=2em]
        \item \textbf{Implement the type-2 flag SDP for small $m$.}  Start with $m=7,9,11,13$, where independent partial or complete arguments exist, and force the KKT square $\mathcal S_m$ into the ansatz.

        \item \textbf{Study vanishing constraints at Turan graphons.}  The certificate must vanish at every balanced $T_r$.  This should dramatically reduce the search space.

        \item \textbf{Look for a symbolic pattern.}  Compare SDP certificates across $m$.  In particular, test whether the same combination of $\mathcal S_m$ and even path-layer surplus terms appears repeatedly.

        \item \textbf{Prove stability from triangle or clique-density stability in residual regimes.}  For $C_{13}$, the remaining cases appear close to triangle-density extremal graphons.  A stability transfer may close that fixed case.

        \item \textbf{Connect with the complement identity.}  Relate $x_n-c_m$ to type-2 squares and avoid proof strategies that discard it.  The $C_{13}$ diagnostics show that discarding this term is too wasteful.
\end{enumerate}

\begin{thebibliography}{9}

\bibitem{RazborovFlag}
A. A. Razborov,
\newblock Flag algebras,
\newblock \emph{Journal of Symbolic Logic} 72 (2007), no. 4, 1239--1282.

\bibitem{RazborovTriangles}
A. A. Razborov,
\newblock On the minimal density of triangles in graphs,
\newblock \emph{Combinatorics, Probability and Computing} 17 (2008), no. 4, 603--618.

\bibitem{ReiherClique}
C. Reiher,
\newblock The clique density theorem,
\newblock \emph{Annals of Mathematics} 184 (2016), no. 3, 683--707.

\bibitem{KimLee}
J. Kim and J. Lee,
\newblock Extended commonality of paths and cycles via Schur convexity,
\newblock arXiv:2210.00977.

\bibitem{BDLP}
P. Bennett, A. Dudek, B. Lidicky, and O. Pikhurko,
\newblock Minimizing the number of $C_5$'s in graphs with given edge-density,
\newblock arXiv:1803.00165.

\end{thebibliography}

\end{document}
